\documentclass[11pt]{article}
\usepackage[margin=2.3cm]{geometry}
\usepackage{amsmath,amssymb,booktabs}
\allowdisplaybreaks
\input{../build/R07_numbers.tex}
\title{XVI-R R07: the reduced $hhc$ loop residual, the p22 $\kappa^2$ evanescent counterterms,\\
and a second contraction artefact ($\bar H$: 18 $\to$ \HB)}
\author{XVI-R (local landmark; imports XVI-Q frozen at Q11; not published)}
\date{}
\begin{document}\maketitle
\begin{abstract}
R07 addresses two fronts. \textbf{Front 1} is the $h\!\cdot\!h\!\cdot\!c$ loop residual in the reduced variables $\bar\Gamma$.
It is \textbf{CLOSED}. The reduced row is $\bar R=A+B+C+D_{EH}$. The difference between the full and reduced rows is
exactly $D_{gf}+\sigma_F E$, which is purely evanescent (\GFE{} terms) and explains the R06 mismatch. The tree splits
exactly. The inhomogeneous reduced fit, including the ghost equation, is solvable with a purely evanescent counterterm
that is verified on every row.
\textbf{Front 2} is consistency with the p22 $\kappa^2$ evanescent counterterms. It is \textbf{CLOSED} at the level of
three-point rows, one loop and UV poles. In our calculus the p22 operators are exactly minus the $\tilde g$-structure
parts of the two-point poles. For D01, D03, D11, D13 and D06 the ratio to XVI-P equals the frozen Q05 validation ratio.
When inserted into the $\kappa^3$ rows, they admit an evanescent cubic completion.
On the way we found a \textbf{second contraction artefact}. The frozen probe contraction also mistreats the external
$\tilde K\leftrightarrow E$ graviton line. Applying the R05 rule ($\bar\eta\cdot\bar\eta$, $\hat\eta\cdot\hat\eta$,
mixed $=0$) on every external probe line changes the reduced $hhc$ block (rank \BRo$\to$\BRn, joint rank \BRj). The
cohomology becomes $\dim\bar H=\HB$, which \emph{supersedes} the 18 of R05/R06 (disclosure R07-D1). Outcome:
\textbf{Partial}.
\end{abstract}

\section{Scope}
Premises are XVI-Q Q11 (pinned zip) and the R01--R06 layers, whose gates must ACCEPT. OpenAI.tex and the new sheets are
not premises. O.1 is CLOSED only in its scoped sense. Bulk symmetry is not restored and the insertion is frozen. No frozen
code is edited: module attributes are swapped only for the duration of a call. All R07 code is in \texttt{r07/}.

\section{The second contraction artefact (R07-D1, R07-D2)}
\texttt{gravity.contract} rewrites each probe pair by dropping every $T$ copy and treating the $H$ copy as the full
four-dimensional dot. This is exact only when \emph{both} sides carry four-dimensional dots ($H+T$ pairs). R05 found
the failure on the $c^*$ line. The same failure occurs on the external graviton-probe line $\tilde K\leftrightarrow E$
(\texttt{rows.contract\_id(\dots,'K','E','h')}). There the $\tilde K$ side can carry a $\hat g$ contraction or a
$\tilde g$ contraction:
\begin{itemize}
\item every evanescent $X_K$ monomial with $T(a_K,\cdot)$, in the $D_{EH}$ channel of the reduced row and in the F and
KKc $\bar\sigma$-images;
\item the $E_2$ structure of the D13 pole, in channel $C$;
\item the $\hat g$ from symmetric integration.
\end{itemize}
With $\eta=\bar\eta+\hat\eta$ the correct rule is $H\cdot H\to H$, $T\cdot T\to T$, mixed $\to0$. For example,
$T(K,c)\,[H(E,x)+T(E,x)]\to T(c,x)$, whereas the frozen rule gives $0$. This is unit-tested. On
$H+T$ pairs both rules agree, so internal loop lines are unaffected.

\textbf{Effects.} The reduced $hhc$ linear block changes from rank \BRo{} to \BRn{}, and its row space changes (joint rank
\BRj). The ghost-row blocks equal those of R05. The loop channels $C$ and $D_{EH}$ change by evanescent terms only (their
4D projections are unchanged), so every earlier 4D-closure statement stands.

\section{Corrected classification}
We repeat R06's complete-$Y$ computation (evanescent and 4D generators: \NGF{} F, \NGG{} G, \NGK{} KKc with nonzero image;
\PEC{} combinations with purely evanescent images) under both rules.
\begin{center}\begin{tabular}{lccccc}\toprule
rule & $\dim\bar N_g$ & $\dim T$ & blind & $T\cap$blind & $\dim\bar H$\\\midrule
$c^*$ line only (R05/R06) & 133 & 113 & 38 & 36 & 18\\
all external probe lines (R07) & \NB & \TB & \BG & \TCB & \textbf{\HB}\\\bottomrule\end{tabular}\end{center}
The first row reproduces R05/R06 exactly; this is the sanity check. Under the corrected rule $\bar\sigma$ remains nilpotent
on the complete $Y$, since every image lies in $\bar N_g$. R05's 18 representatives are no longer cocycles of the
corrected system. The completeness arguments of R06 (total derivatives, field redefinitions/EOM, the formal $H/T$
algebra, parity) are unchanged, so the corrected statement is: \textbf{the invariant, parity-even evanescent cohomology
at $O(\kappa^3)$ on the complete local $Y$-space has dimension \HB}. ``Physical'' is meant only in this sense; no
direction is certified as an observable. The basis is in \texttt{R07\_H\_basis.json}.

\section{Front 1: the reduced $hhc$ loop residual}
With $\tilde K=K+\bar cF'(h)$ the gauge-fixing part of $S_2$ and the $\bar c$ channel are absorbed into $\tilde K$:
$$\bar R=A+B_2+B_3+C+D_{2,EH}+D_{3,EH}.$$
\begin{itemize}
\item Tree: $A_t+B_t^{EH}=0$ and $B_t^{gf}+\sigma_FE_t=0$ hold separately and exactly.
\item Loop: full $-\bar R=D_{gf}+\sigma_F E$ exactly (\GFE{} terms, all evanescent). R06 had paired the full residual with
the reduced channels; that mismatch is now explained.
\item Fit with the $c^*$-line rule: $\bar R$ (\RBo{} terms) plus reduced channels, both ghost rows, Bose symmetry and the
ghost equation is solvable, with $\dim\bar N_g=133$.
\item Fit with the corrected rule: $\bar R$ (\RBn{} terms, 4D part zero) is solvable with the ghost equation, and
$\dim\bar N_g=\NB$. The particular counterterm has \FKa/\FCa/0/\FHa{} terms in $K/C/S/H$, is purely evanescent, and
makes every row vanish.
\end{itemize}
\textbf{Outcome: CLOSED.}

\section{Front 2: p22 consistency}
\subsection{Identification}
In our calculus the $\tilde g$-structure part of a pole $o$ is $P_{ev}(o)=o-\mathrm{lift}(o|_{T=0})$, where lift sends
$H(x,y)\mapsto H+T$ on polarisation and probe pairs, because the 4D structures are $g_4$ and $p$ monomials.
\begin{center}\begin{tabular}{lccc}\toprule
graph & own/XVI-P (full) & own/XVI-P (evanescent part) & p22 CT $=-$pole\\\midrule
D01 graviton & $i$ & $i$ & yes\\
D03 ghost loop & $i$ & $i$ & yes\\
D11 scalar & $i$ & $i$ & yes\\
D13 $\tilde K c$ ($E_2$) & $-i$ & $-i$ & yes\\
D06 $\bar cc$ (SE) & $i$ & $i$ & yes\\\bottomrule\end{tabular}\end{center}
The evanescent ratios equal the frozen Q05 normalisation ratios. The p22 counterterms are therefore exactly
$X_2=-P_{ev}(\text{pole})$ in our normalisation.

\subsection{Cubic completion}
We insert $X_2$ in place of every two-point pole in the $\kappa^3$ rows: $hhc$ channels $B$ and $C$, and ghost-row
channels $K$ and $c_2$. All genuine three-point loops are set to zero and the corrected rule is used. The inhomogeneity
$\Xi$ has \XIh/\XIk/\XIb{} terms ($hhc$/$ccK$/$cc\bar c$), all evanescent. The system $\bar\sigma X_3'+\Xi=0$, with
Bose symmetry and the ghost equation, is \textbf{solvable}. The completion $X_3'$ (\PKa/\PCa/0/\PHa{} terms) is purely
evanescent and verified. It is unique modulo $\bar N_g$; its class in $\bar H$ (dimension \HB) is not fixed by p22.
\textbf{Outcome: CLOSED} at three-point, one-loop pole level. Not implied: four-point rows, finite parts, two loops.

\section{Disclosures}
\begin{description}
\item[R07-D1] The hat rule is extended to every external probe line. The reduced $hhc$ block and the F/KKc images are
rebuilt, and $\bar H$ goes from 18 to \HB. The R05/R06 count of 18 is superseded. Frozen code is not edited.
\item[R07-D2] Loop channels $C$ and $D_{EH}$ change (evanescent only; 4D closure unaffected).
\item[R07-D3] $\bar R$ is defined as $A+B+C+D_{EH}$; full $-$ reduced $=D_{gf}+\sigma_FE$.
\item[R07-D4] All R07 code is in \texttt{r07/}; the R01--R06 gates are untouched.
\end{description}

\section{Open items}
Parity-odd evanescent cohomology; counting in a Gram-reduced scheme; four-point rows; finite parts. Bulk symmetry stays
not restored, the insertion stays frozen, and Route B is not discharged.
\end{document}
